\begin{thebibliography}{99}\small
\bibitem{prat2017} Prat Y, Taub M, Pratt E, Yovel Y. An annotated dataset of Egyptian fruit bat vocalizations across varying contexts and during vocal ontogeny. \emph{Scientific Data}. 2017;4:170143. doi:10.1038/sdata.2017.143
\bibitem{prat2016} Prat Y, Taub M, Yovel Y. Everyday bat vocalizations contain information about emitter, addressee, context, and behavior. \emph{Scientific Reports}. 2016;6:39419. doi:10.1038/srep39419
\bibitem{elie2024} Elie JE, Muroy SE, Genzel D, et al. Role of auditory feedback for vocal production learning in the Egyptian fruit bat. \emph{Current Biology}. 2024;34(17):4062--4070.e7. doi:10.1016/j.cub.2024.07.053 (data: doi:10.6084/m9.figshare.25886023).
\bibitem{aves2023} Hagiwara M. AVES: Animal Vocalization Encoder Based on Self-Supervision. In: \emph{ICASSP 2023 -- IEEE International Conference on Acoustics, Speech and Signal Processing}. 2023:1--5. doi:10.1109/ICASSP49357.2023.10095642; arXiv:2210.14493.
\bibitem{bezdek1984} Bezdek JC, Ehrlich R, Full W. FCM: The fuzzy c-means clustering algorithm. \emph{Computers \& Geosciences}. 1984;10(2--3):191--203.
\bibitem{xiebeni1991} Xie XL, Beni G. A validity measure for fuzzy clustering. \emph{IEEE Trans Pattern Anal Mach Intell}. 1991;13(8):841--847.
\bibitem{rousseeuw1987} Rousseeuw PJ. Silhouettes: a graphical aid to the interpretation and validation of cluster analysis. \emph{J Comput Appl Math}. 1987;20:53--65.
\bibitem{hubert1985} Hubert L, Arabie P. Comparing partitions. \emph{J Classification}. 1985;2:193--218.
\bibitem{mcfee2015} McFee B, et al. librosa: Audio and music signal analysis in Python. \emph{Proc 14th Python in Science Conf}. 2015:18--25.
\bibitem{pedregosa2011} Pedregosa F, et al. Scikit-learn: Machine learning in Python. \emph{J Mach Learn Res}. 2011;12:2825--2830.
\bibitem{tibshirani2005} Tibshirani R, Walther G. Cluster validation by prediction strength. \emph{J Comput Graph Stat}. 2005;14(3):511--528.
\bibitem{hennig2007} Hennig C. Cluster-wise assessment of cluster stability. \emph{Comput Stat Data Anal}. 2007;52(1):258--271.
\bibitem{hartigan1985} Hartigan JA, Hartigan PM. The dip test of unimodality. \emph{Ann Statist}. 1985;13(1):70--84.
\bibitem{bergsma2013} Bergsma W. A bias-correction for Cram\'er's V and Tschuprow's T. \emph{J Korean Stat Soc}. 2013;42:323--328.
\bibitem{tibshirani2001} Tibshirani R, Walther G, Hastie T. Estimating the number of clusters in a data set via the gap statistic. \emph{J R Stat Soc B}. 2001;63(2):411--423.
\bibitem{campello2013} Campello RJGB, Moulavi D, Sander J. Density-based clustering based on hierarchical density estimates. \emph{PAKDD 2013, LNCS}. 2013;7819:160--172. doi:10.1007/978-3-642-37456-2\_14
\bibitem{seabold2010} Seabold S, Perktold J. statsmodels: Econometric and statistical modeling with Python. \emph{Proc 9th Python in Science Conf}. 2010. doi:10.25080/Majora-92bf1922-011
\bibitem{wadewitz2015} Wadewitz P, Hammerschmidt K, Battaglia D, Witt A, Wolf F, Fischer J. Characterizing vocal repertoires --- hard vs.\ soft classification approaches. \emph{PLoS ONE}. 2015;10(4):e0125785. doi:10.1371/journal.pone.0125785
\bibitem{fischer2017} Fischer J, Wadewitz P, Hammerschmidt K. Structural variability and communicative complexity in acoustic communication. \emph{Animal Behaviour}. 2017;134:229--237. doi:10.1016/j.anbehav.2016.06.012
\bibitem{sainburg2020} Sainburg T, Thielk M, Gentner TQ. Finding, visualizing, and quantifying latent structure across diverse animal vocal repertoires. \emph{PLoS Comput Biol}. 2020;16(10):e1008228. doi:10.1371/journal.pcbi.1008228
\bibitem{thomas2022} Thomas M, Jensen FH, Averly B, et al. A practical guide for generating unsupervised, spectrogram-based latent space representations of animal vocalizations. \emph{J Anim Ecol}. 2022;91(8):1567--1581. doi:10.1111/1365-2656.13754
\bibitem{goffinet2021} Goffinet J, Brudner S, Mooney R, Pearson J. Low-dimensional learned feature spaces quantify individual and group differences in vocal repertoires. \emph{eLife}. 2021;10:e67855. doi:10.7554/eLife.67855
\bibitem{best2023} Best P, Paris S, Glotin H, Marxer R. Deep audio embeddings for vocalisation clustering. \emph{PLoS ONE}. 2023;18(7):e0283396. doi:10.1371/journal.pone.0283396
\bibitem{stowell2022} Stowell D. Computational bioacoustics with deep learning: a review and roadmap. \emph{PeerJ}. 2022;10:e13152. doi:10.7717/peerj.13152
\bibitem{adolfsson2019} Adolfsson A, Ackerman M, Brownstein NC. To cluster, or not to cluster: an analysis of clusterability methods. \emph{Pattern Recognition}. 2019;88:13--26. doi:10.1016/j.patcog.2018.10.026
\bibitem{liu2008} Liu Y, Hayes DN, Nobel A, Marron JS. Statistical significance of clustering for high-dimension, low-sample size data. \emph{J Am Stat Assoc}. 2008;103(483):1281--1293.
\bibitem{chari2023} Chari T, Pachter L. The specious art of single-cell genomics. \emph{PLoS Comput Biol}. 2023;19(8):e1011288. doi:10.1371/journal.pcbi.1011288
\bibitem{roberts2017} Roberts DR, et al. Cross-validation strategies for data with temporal, spatial, hierarchical, or phylogenetic structure. \emph{Ecography}. 2017;40(8):913--929. doi:10.1111/ecog.02881
\bibitem{hurlbert1984} Hurlbert SH. Pseudoreplication and the design of ecological field experiments. \emph{Ecol Monogr}. 1984;54(2):187--211.
\bibitem{mcinnes2018} McInnes L, Healy J, Melville J. UMAP: Uniform Manifold Approximation and Projection for dimension reduction. arXiv:1802.03426, 2018.
\end{thebibliography}
